\documentclass[11pt]{article}
\usepackage[margin=1in]{geometry}
\usepackage{amsmath,amssymb,amsthm,mathtools}
\usepackage{graphicx}
\usepackage{microtype}
\usepackage[colorlinks=true,linkcolor=blue,urlcolor=blue,citecolor=blue]{hyperref}

\newtheorem{theorem}{Theorem}
\newtheorem{lemma}[theorem]{Lemma}
\newtheorem{remark}[theorem]{Remark}
\newcommand{\R}{\mathbb R}
\newcommand{\opnorm}[3]{\lVert #1\rVert_{#2\to #3}}

\title{Bounded $\Gamma_n$ forces bounded $\Delta_n$:\\
a negative answer to Remark 5.5 of arXiv:math/0610427}
\author{Candidate full disproof of the conjectured existence statement}
\date{17 August 2026}

\begin{document}
\maketitle

\noindent\fbox{\parbox{0.94\linewidth}{\textbf{Status: candidate full
disproof, high confidence.}  For every probability measure $P$ considered in
the source paper,
\[
 \lVert\Delta_n(P)\rVert_\infty
 =\opnorm{\Gamma_n(P)}{2}{\infty}^{,2}
 \leq \lVert\Gamma_n(P)\rVert_2^2.
\]
Hence a uniformly bounded sequence
$\lVert\Gamma_n(P_n)\rVert_2$ can never coexist with an unbounded sequence
$\lVert\Delta_n(P_n)\rVert_\infty$.  The sequence conjectured in Remark 5.5
does not exist.}}

\section{\hspace{0.25em}The source question}

Kontorovich associates to a probability measure $P$ on a product space an
upper-triangular matrix $\Delta_n(P)$ whose diagonal entries are one and whose
strictly upper-triangular entries are mixing coefficients
$\bar\eta_{ij}\in[0,1]$.  Thus
\begin{equation}
 \lVert\Delta_n(P)\rVert_\infty
 =\max_{1\leq i<n}\left(1+\sum_{j=i+1}^n\bar\eta_{ij}\right).
 \label{eq:delta}
\end{equation}
The paper also uses Samson's matrix $\Gamma_n(P)$, defined by taking the
entrywise nonnegative square root,
\begin{equation}
       (\Gamma_n(P))_{ij}=\sqrt{(\Delta_n(P))_{ij}},
 \label{eq:gamma}
\end{equation}
and writes $\lVert\Gamma_n(P)\rVert_2$ for its induced Euclidean (spectral)
operator norm.

Remark 5.5 asks whether there is a sequence $(P_n)$ such that
$\lVert\Gamma_n(P_n)\rVert_2$ is bounded while
$\lVert\Delta_n(P_n)\rVert_\infty$ is unbounded, and conjectures that such a
sequence exists \cite[p.~11, Remark 5.5]{Kontorovich}.  The latest arXiv
version is v2, dated 14 October 2006, and retains the question.

\begin{center}
\includegraphics[width=0.96\linewidth]{figures/open_problem_crop.png}
\end{center}

\section{\hspace{0.25em}The matrix identity}

The obstruction does not require the probabilistic structure of the entries.

\begin{lemma}\label{lem:matrix}
Let $D=(d_{ij})$ be an $m\times n$ matrix with nonnegative entries, and set
$G=(\sqrt{d_{ij}})$ entrywise.  Then
\begin{equation}
                 \lVert D\rVert_\infty
             =\opnorm{G}{2}{\infty}^{,2}
             \leq \opnorm{G}{2}{2}^{,2}.
 \label{eq:master}
\end{equation}
Here $\lVert D\rVert_\infty$ is the induced $\ell_\infty$ operator norm.
\end{lemma}

\begin{proof}
For any real matrix $A$,
\begin{equation}
       \opnorm{A}{2}{\infty}
       =\max_{1\leq i\leq m}\lVert \operatorname{row}_i(A)\rVert_2.
 \label{eq:row}
\end{equation}
Indeed, Cauchy--Schwarz gives the upper bound in \eqref{eq:row}; equality is
obtained by taking the unit vector parallel to a row of maximum Euclidean
norm.  Applying this to $G$ gives
\[
 \opnorm{G}{2}{\infty}^{,2}
   =\max_i\sum_j |g_{ij}|^2
   =\max_i\sum_j d_{ij}
   =\lVert D\rVert_\infty,
\]
where the last equality uses $d_{ij}\geq0$.  Finally,
$\lVert y\rVert_\infty\leq\lVert y\rVert_2$ for every $y\in\R^m$, so
\[
  \opnorm{G}{2}{\infty}
   =\sup_{\lVert x\rVert_2=1}\lVert Gx\rVert_\infty
   \leq\sup_{\lVert x\rVert_2=1}\lVert Gx\rVert_2
   =\opnorm{G}{2}{2}.
\]
Squaring proves \eqref{eq:master}.
\end{proof}

\section{\hspace{0.25em}Full negative answer}

\begin{theorem}\label{thm:answer}
For every $n$ and every probability measure $P$ for which the source
matrices are defined,
\begin{equation}
 \boxed{\quad
 \lVert\Delta_n(P)\rVert_\infty
 =\opnorm{\Gamma_n(P)}{2}{\infty}^{,2}
 \leq\lVert\Gamma_n(P)\rVert_2^2.
 \quad}
 \label{eq:answer}
\end{equation}
Consequently, if
$\sup_n\lVert\Gamma_n(P_n)\rVert_2\leq C<\infty$, then
$\sup_n\lVert\Delta_n(P_n)\rVert_\infty\leq C^2$.  In particular, the
sequence requested and conjectured in Remark 5.5 does not exist.
\end{theorem}

\begin{proof}
Apply Lemma~\ref{lem:matrix} with $D=\Delta_n(P)$ and
$G=\Gamma_n(P)$, using the entrywise relation \eqref{eq:gamma}.  Formula
\eqref{eq:delta} omits the last row from its displayed maximum.  That row has
sum one, while every earlier row also contains its diagonal one, so including
the last row does not change the maximum.  Thus Lemma~\ref{lem:matrix} applies
without any convention mismatch.  The uniform implication follows directly
from \eqref{eq:answer}.
\end{proof}

\section{\hspace{0.25em}Checks, sharpness, and relation to Theorem 5.3}

\paragraph{Sharpness.}
For a product measure, $\Delta_n=\Gamma_n=I_n$.  Equality holds throughout
\eqref{eq:answer}; the constant one cannot be improved.

\paragraph{No hidden measure-theoretic assumption.}
Essential suprema enter only in defining the numbers $\bar\eta_{ij}$.  Once
the finite matrices have been formed, the proof is an exact algebraic
identity.  It applies equally to finite, countable, or continuous underlying
state spaces whenever the source definitions apply.

\paragraph{Compatibility with Theorem 5.3.}
The source proves that the ratio
\[
       R_n=\frac{\lVert\Gamma_n(P_n)\rVert_2}
                  {\lVert\Delta_n(P_n)\rVert_\infty}
\]
can tend either to zero or to infinity.  This does not conflict with
\eqref{eq:answer}.  The zero-ratio construction has both norms unbounded, with
$\Gamma_n$ permitted to grow on the scale of the square root of $\Delta_n$;
the infinite-ratio construction keeps $\Delta_n$ bounded while $\Gamma_n$
grows.  The missing fourth possibility---bounded $\Gamma_n$ and unbounded
$\Delta_n$---is exactly what \eqref{eq:answer} forbids.

\section{\hspace{0.25em}Literature and version audit}

Exact-phrase searches for the question, together with searches combining the
author, the two matrix symbols, mixing coefficients, and operator norms, found
no explicit published resolution as of 17 August 2026.  The author's 2007
thesis recalls the $\Gamma_n$--$\Delta_n$ comparison and Theorem 5.3 in
Section 5.7, but omits this conjecture and does not state
\eqref{eq:answer} as an answer \cite[\S5.7]{Thesis}.  The related later paper
on prescribed $\eta$-mixing coefficients addresses a different realization
problem \cite{Prescribed}.  This packet therefore records
Theorem~\ref{thm:answer} as an apparently new explicit negative answer,
subject to independent literature review.

\begin{thebibliography}{9}
\bibitem{Kontorovich}
L.~Kontorovich,
\emph{Metric and Mixing Sufficient Conditions for Concentration of Measure},
arXiv:math/0610427v2 (2006),
\url{https://arxiv.org/abs/math/0610427}.

\bibitem{Thesis}
L.~Kontorovich,
\emph{Measure Concentration of Strongly Mixing Processes with Applications},
Ph.D. thesis, Carnegie Mellon University (2007),
\url{https://www.cs.bgu.ac.il/~karyeh/thesis.pdf}.

\bibitem{Prescribed}
L.~Kontorovich,
\emph{Constructing processes with prescribed mixing coefficients},
arXiv:0711.0986 (2007),
\url{https://arxiv.org/abs/0711.0986}.
\end{thebibliography}

\end{document}
